\fbox{例題 1.6}

データを大きさ順に並べると
\begin{equation*}
  5.9, 6.8, 7.0, 7.2, 7.4, 7.5, 7.8, 7.8, 8.1, 8.3
\end{equation*}
\begin{equation*}
  8.4, 8.5, 8.8, 8.9, 9.0, 9.2, 9.3, 9.8, 10.1, 10.8
\end{equation*}

\begin{equation*}
  \text{トリム平均} = \frac{6.8 + 7.0 + \cdots + 10.1}{20 - 2} = 5.66
\end{equation*}
\begin{equation*}
\text{メディアン}Me = \frac{8.3 + 8.4}{2} = 8.35
\end{equation*}
\begin{equation*}
  \text{第1四分位点} = \frac{7.4 + 7.5}{2} = 7.45
\end{equation*}
\begin{equation*}
  \text{第3四分位点} = \frac{9.0 + 9.2}{2} = 9.10
\end{equation*}
\begin{equation*}
  \text{四分位範囲} = 9.10 - 7.45 = 1.65
\end{equation*}
